\documentclass[../main]{subfiles}
\begin{document}
\chapter{Mantenimiento}
\img{images/07/mantenimiento.jpg}{Mantenimiento predictivo: medición de vibraciones en acople rígido.}
\info{Mantenimiento}{Conservación de una cosa en buen estado o en una situación determinada para evitar su degradación.}

\section{Tipos de fallas}
\info{Falla}{Es la incapacidad del activo de cumplir con su función, cualquier función. Son todos los estados indeseables del sistema o equipo}

\subsection{Falla técnica}
Una falla técnica a diferencia de la falla funcional permite que la máquina siga funcionando, pero con algunas alteraciones que son perceptibles, ruido, vibraciones, y que en definitiva van a producir un fallo funcional en algún momento.

\subsection{Falla funcional}
Deja de cumplir con la función, sin embargo no necesariamente significa que el activo deja de producir. Aquí es muy importante definir los requerimientos de desempeño. Son resultado de la incapacidad de los equipos para dar su nivel de desempeño esperado. Son un síntoma o evento observable por el personal de producción y mantenimiento.

\section{Tipos de mantenimiento}

\subsection{Mantenimiento correctivo}

Se denomina mantenimiento correctivo a aquel que corrige los defectos observados en los equipamientos o instalaciones, es la forma más básica de mantenimiento y consiste en localizar averías o defectos para corregirlos o repararlos. Cuando hay un fallo funcional se repara inmediatamente la máquina.

\subsection{Mantenimiento preventivo}
En las operaciones de mantenimiento el mantenimiento preventivo es el destinado a la conservación de equipos o instalaciones mediante la realización de revisión y limpieza que garanticen su buen funcionamiento y fiabilidad. El mantenimiento preventivo se realiza en equipos en condiciones de funcionamiento, por oposición al mantenimiento correctivo que repara o pone en condiciones de funcionamiento aquellos que dejaron de funcionar o están dañados.

\subsection{Mantenimiento predictivo}
Son una serie de acciones que se toman, y técnicas que se aplican, con el objetivo de detectar posibles fallos y defectos de maquinaria en las etapas incipientes, para evitar que estos fallos se manifiesten en uno más grande durante su funcionamiento, evitando que ocasionen paros de emergencia y tiempos muertos, causando impacto financiero negativo. Su misión es conservar un nivel de servicio determinado en los equipos programando las revisiones en el momento más oportuno. Suele tener un carácter sistemático, es decir, se interviene aunque el equipo no haya dado ningún síntoma de tener problemas.

\subsection{Mantenimiento legal}

Se denomina mantenimiento legal, al tipo de mantenimiento preventivo obligatorio que las legislaciones sobre seguridad de equipos e instalaciones industriales obligan que deben realizarse de forma periódica por parte de empresas o personal autorizado ajenas a la empresa propiedad de las instalaciones o equipos. 

\subsection{Mantenimiento subcontratado}
También llamado subcontratación o externalización, es un proceso económico empresarial en el cual una empresa delega ciertas tareas, no propias de su función pero igualmente importantes e indispensables, a una empresa externa especializada en la prestación de esos servicios.

\actividad{ Actividad para el proyecto.}
\begin{enumerate}
    \item Enumerar las partes de su proyecto que puedan ser susceptibles a producir fallos funcionales y seleccione el tipo de mantenimiento apropiado, justificar su respuesta.
    \item Elaborar un código QR para los datos básicos de la máquina.
    \item Realizar una ficha de equipo para todas las máquinas que estén presente en su proyecto. Deberá contener:
    \begin{enumerate}
        \item Codificación de equipo.
        \item Foto del equipo o dibujo.
        \item Peso, y dimensiones.
        \item Dirección y datos del vendedor y/o fabricante.
        \item Insumos y elementos necesarios para su funcionamiento.
        \item códigos de repuestos.
        \item Ejemplo de Fija de equipo a realizar:
        \begin{center}
        \includegraphics[]{images/07/fichamant.png}
        \end{center}
        
    \end{enumerate}
\end{enumerate}

\actividad { Cuestionario: investigar}
\begin{enumerate}
    \item  ¿Qué es la obsolescencia programada?
    \item ¿Qué es la obsolescencia tecnológica?
    \item ¿Para que se utiliza el analisis de criticidad en mantenimiento?
    \item ¿Qué es una ficha de equipo?
    \item ¿Qué es la burocracia?
    \item ¿Porqué es importante tener una ficha de equipo de mantenimiento?
    \item ¿Qué es la codificiación de equipos?
    \item ¿Qué es el código de barra, QR cómo funcionan y para que puede utilizarse?
    \item ¿Qué es el mantenimiento total?
    
    \end{enumerate}


\end{document}